% ======================================================================
% Ontologie der Kontinua -- Core 1.3
% Cross-K-Modul: crossk_k2_k3.tex
% Cross-Level-Relationen zwischen K2 und K3
% ======================================================================

\subsubsection{K2–K3-Querverbindung}
\label{sec:crossk-k2-k3}

Der Übergang $K_2 \rightarrow K_3$ beschreibt die Entstehung des chemischen
Kontinuums aus einem physischen Kontinuum, das bereits Geometrie, Felder und
dynamische Flüsse besitzt. Während $K_2$ physikalische Freiheitsgrade enthält
(räumliche Achsen, energetische Potentiale, Feldgradienten, Perkolationsstruktur),
führt $K_3$ ein:
\begin{itemize}
  \item chemische Zusammensetzungsräume,
  \item Bindungsenergien und molekulare Potentiale,
  \item Reaktionsschwellen und stöchiometrische Nebenbedingungen,
  \item stabile Cluster und Netzwerkbindungen,
  \item neue Flussklassen (Reaktion, Diffusions-Reaktion, Ladungsübertragung),
  \item einen erweiterten Raum $\Omega(K_3)$, strukturiert durch chemische
        Konnektivität.
\end{itemize}

Der Übergang K₂→K₃ ist die erste große Steigerung struktureller Komplexität:
Physikalische Felder werden zum Substrat chemischer Organisation.

% ----------------------------------------------------------------------
\subsubsection{1. Stufen und ihre Rollen}

\paragraph{K2 (physikalisches Kontinuum).}
$K_2$ liefert:
\begin{itemize}
  \item mehrdimensionale Achsen (geometrische Koordinaten),
  \item energetische Potentiale $P_2$,
  \item Flüsse $J_2$ (Transport, Diffusion, Feldausbreitung),
  \item Perkolationsstruktur der Konnektivität,
  \item kausale Ordnung von Ereignissen.
\end{itemize}

\paragraph{K3 (chemisches Kontinuum).}
$K_3$ baut auf $K_2$ auf und fügt hinzu:
\begin{itemize}
  \item neue Achsen, die chemische Spezies, Zusammensetzung und Ladung
        beschreiben,
  \item Potentiale für Bindungsenergie und Reaktionslandschaften,
  \item Schwellen $\Theta_3$ für Reaktionsdurchführbarkeit und molekulare Stabilität,
  \item Flüsse $J_3$ für Reaktionen, Dissoziation und Ladungsübertragung,
  \item stabile Cluster und Reaktionszyklen als chemische Pfade.
\end{itemize}

$K_3$ ist das erste Kontinuum, das Templating, Katalyse und
informationsähnliche Persistenz über molekulare Zustände unterstützt.

% ----------------------------------------------------------------------
\subsubsection{2. Gemeinsame und vererbte Achsen sowie Schwellen}

\paragraph{Achsen.}
Die räumlichen Achsen von $K_2$ werden vollständig geerbt:
\[
A_{\mathrm{geom}}(K_2) \subset A(K_3).
\]
$K_3$ ergänzt chemische Achsen wie:
\begin{itemize}
  \item Zusammensetzung/Stöchiometrie,
  \item Ladungszustand,
  \item Bindungskonfiguration,
  \item Reaktionskoordinate(n).
\end{itemize}

\paragraph{Schwellen.}
Chemische Schwellen $\Theta_3$ erweitern physikalische Schwellen $\Theta_2$.
Wesentliche vererbte Nebenbedingungen sind:
\begin{itemize}
  \item energetische Schwellen: unzureichendes $P_2$ verhindert Bindungsbildung,
  \item Perkolationsschwellen: Konnektivität von Interaktionsregionen muss den
        kritischen Wert $p_c$ überschreiten,
  \item Gradientschwellen: ausreichender Fluss oder Kollisionsfrequenz ist für
        Reaktionsaktivierung erforderlich.
\end{itemize}

Scheitert eine physikalische Schwelle, kann keine chemische Struktur entstehen:
\[
\Theta_2^{\mathrm{death}} \Rightarrow \Omega(K_3)=\varnothing.
\]

% ----------------------------------------------------------------------
\subsubsection{3. Cross-Level-Flüsse, Zyklen und Spannungen}

\paragraph{Flüsse.}
Flüsse in $K_3$ basieren auf $J_2$, enthalten nun aber:
\begin{itemize}
  \item Reaktionsflüsse entlang chemischer Koordinaten,
  \item Diffusions-Reaktions-Kopplung,
  \item Ladungsfluss und Redox-Flüsse,
  \item Flüsse zwischen Potentialtälern in der Reaktionslandschaft.
\end{itemize}

Formal:
\[
J_3 = (J_2,\ J_{\mathrm{react}},\ J_{\mathrm{charge}},\ J_{\mathrm{bond}}).
\]

\paragraph{Zyklen.}
Chemische Zyklen sind höherdimensional als physikalische Schleifen. Sie
representieren:
\[
C_3 : \text{reaktive Sequenz } A \rightarrow B \rightarrow C \rightarrow A.
\]
Chemische Zyklen erfordern $K_2$-Transport plus $K_3$-Reaktionspfade.

Das Dasein eines stabilen Zyklus in $K_3$ hängt ab von:
\[
\oint_{\gamma} J_3 \cdot dA > 0
\quad\text{und}\quad
\text{alle Zwischenzustände erfüllen }\Theta_3.
\]

\paragraph{Spannung.}
Die Cross-Level-Spannung $T_{2,3}$ misst die Kompatibilität zwischen physikalischen
Feldern und chemischen Nebenbedingungen:
\[
T_{2,3} = f(P_3 - P_2,\ \Theta_3 - \Theta_2,\ A_{\mathrm{chem}}).
\]
Überschreitet $T_{2,3}$ die Schwelle, kollabiert das chemische Kontinuum zurück
in ein rein physikalisches Regime (keine stabilen Moleküle oder Reaktionen).

% ----------------------------------------------------------------------
\subsubsection{4. Geburts- und Todesbedingungen über die Stufen}

\paragraph{Geburt von \texorpdfstring{$K_3$}{K_3}.}
Chemische Struktur entsteht, wenn:
\[
T_2 > \Theta_{2,\mathrm{dim}}
\quad \text{und} \quad
A_{\mathrm{chem}} \in M_3 \setminus A(K_2).
\]
Das entspricht dem Auftreten neuer Freiheitsgrade im Zusammenhang mit
Bindungsenergien und Zusammensetzung.

Der Zustandsraum erweitert sich:
\[
\Omega(K_3) = \Omega(K_2) \cup \Delta\Omega_{\mathrm{chem}}.
\]

\paragraph{Todesausbreitung.}
Wenn $\Theta_2$ versagt, ist chemische Organisation unmöglich.
Wenn $\Theta_3$ versagt (z.\,B. wegen unzureichender Bindungsenergie,
inkompatibler Felder oder exzessivem thermischen Rauschen), kollabiert nur
$K_3$, während $K_2$ unbeeinträchtigt bleibt.

% ----------------------------------------------------------------------
\subsubsection{5. Konzeptionelle Beispiele}

\paragraph{Perkolation ermöglicht chemische Konnektivität.}
Nur wenn Regionen physikalischer Interaktion perkolieren ($p>p_c$), können chemische
Cluster entstehen; andernfalls können sich keine Moleküle stabilisieren.

\paragraph{Reaktionslandschaften.}
Der Übergang führt neue Potentiale:
\[
P_3 = P_2 + E_{\mathrm{bond}} + E_{\mathrm{charge}},
\]
ein, die Täler und Aktivierungsbarrieren ausbilden und Reaktionspfade definieren.

\paragraph{Operatorbild (Erzeugung/Annihilation von Clustern).}
Die Entstehung oder Dissoziation eines Minimalclusters $q_{\mathrm{min}}$ kann über
Operatoren geschrieben werden als:
\[
a^{\dagger} |k\rangle = \sqrt{f(k)+1}\,|k+k_{\mathrm{unit}}\rangle,
\qquad
a |k\rangle = \sqrt{f(k)}\,|k-k_{\mathrm{unit}}\rangle,
\]
was die quantisierte Änderung der Konnektivität in $\Omega(K_3)$ widerspiegelt.

Diese Beispiele zeigen das Wesen des K2→K3-Übergangs in komprimierter Form:
stabile Cluster, Bindungsenergien und Reaktionspfade entstehen aus physikalischen
Feldern und Geometrie.
